\section{Datos del Alumno, Entidad y Entorno de Practicas}

\subsection{Datos personales del alumno}
\begin{itemize}[leftmargin=1.5em]
    \item Nombre y apellidos: \textbf{[COMPLETAR]}
    \item DNI/NIE: \textbf{[COMPLETAR]}
    \item Titulacion: \textbf{[COMPLETAR]}
    \item Curso: \textbf{[COMPLETAR]}
    \item Correo universitario: \textbf{[COMPLETAR]}
    \item Telefono de contacto: \textbf{[COMPLETAR]}
\end{itemize}

\subsection{Entidad colaboradora y ubicacion}
\begin{itemize}[leftmargin=1.5em]
    \item Entidad: \textbf{ActLab}
    \item Proyecto asignado: \textbf{Didascalia VC - ACTLAB}
    \item Ubicacion del centro de practicas: \textbf{[COMPLETAR DIRECCION / CIUDAD]}
    \item Modalidad: \textbf{[COMPLETAR: presencial / hibrida / remota]}
    \item Periodo de realizacion: \textbf{[COMPLETAR FECHA INICIO - FECHA FIN]}
\end{itemize}

\subsection{Descripcion de la entidad y del departamento asignado}
ActLab es una entidad centrada en el desarrollo de soluciones tecnologicas y experiencias inmersivas, con especial interes en aplicaciones de Realidad Virtual con utilidad formativa. Dentro de este marco, el proyecto \textit{Didascalia VC} se orienta a crear una herramienta de simulacion aplicada a escenarios de aula, donde se representan interacciones y conflictos para apoyar procesos de entrenamiento y toma de decisiones docentes.

En la practica, la participacion se ha realizado en el area tecnica vinculada al desarrollo Unity/VR, con colaboracion transversal entre perfiles de programacion, integracion de contenido (animacion y prefabs), y pruebas funcionales. El rol principal desempe\~nado ha sido de soporte y desarrollo tecnico en arquitectura de sistemas de comportamiento, conexion con backend web y validacion de ejecucion en dispositivo objetivo de VR.

\subsection{Tutorizacion y seguimiento}
\begin{itemize}[leftmargin=1.5em]
    \item Tutor en la empresa: \textbf{[COMPLETAR NOMBRE Y CARGO]}
    \item Tutor academico: \textbf{[COMPLETAR NOMBRE Y DEPARTAMENTO]}
    \item Mecanismo de seguimiento: reuniones periodicas de estado, revisiones tecnicas de tareas y validaciones sobre builds funcionales.
\end{itemize}